\section{Threats to validity}
\label{sec:limitations}

\paragraph{One application.} All live results come from one demo agent with one system prompt and one
tool set. Real assistants differ in prompts, tools, guardrails and in how much telemetry they emit;
the behavioural rules in particular assume fields (provenance, risk class, approval state) that many
applications do not log today.

\paragraph{One frontier provider.} The frontier results are about Claude models behind the Anthropic
API, where a provider classifier sits between the model and the application. On Opus~5.5 most of the
refusal signal is that classifier, so the result describes the deployed system, not model capability
in isolation. Other providers log refusals differently or not at all.

\paragraph{Small cells and fixed attacks.} Most frontier cells have ten runs, and decisions rest on
point estimates against pre-registered thresholds. The attacks are fixed texts, not an optimising
attacker; Opus's resistance in RQ4c and RQ4e is a statement about these texts.

\paragraph{Oracles and labels.} Compliance and harm oracles capture the attacker goal of each scenario,
not every possible harm. RQ4d reports labels from reading all 150 answers; the pre-registered labeller,
a fixed phrase list, gave 46 rather than 47 Haiku relays and mislabelled ten warnings, and \llmd{026}
shares that list.

\paragraph{Dataset labels.} Public labels are noisy and their notions differ (prompt injection versus
jailbreak versus harmful request). Benign corpora are not production traffic.

\paragraph{Author bias.} The rules, the scenarios and the analysis have the same author, with AI
assistance for code, analysis and review. RQ2 and RQ3 rely on third-party attacks; RQ4 Part A
generates attacks automatically; the adversarial review and the pre-registrations are the main
controls, and all of them are in the public repository.
